% ============================================================================
%  math-environments-SN.tex
%
%  SNmono definisce già (in inglese):
%    theorem, lemma, corollary, proposition, example, exercise,
%    remark, note, problem, solution, case, conjecture, question, property
%
%  Qui definiamo solo gli ambienti italiani NON presenti in SNmono.
% ============================================================================

% ── Teorema, Corollario, Proposizione ────────────────────────────────────
% SNmono ha già theorem/corollary/proposition in inglese.
% Definiamo le versioni italiane con contatore condiviso.
\newtheorem{teorema}{Teorema}[section]
\newtheorem{corollario}[teorema]{Corollario}
\newtheorem{proposizione}[teorema]{Proposizione}
\newtheorem{definizione}[teorema]{Definizione}

% ── Osservazione ─────────────────────────────────────────────────────────
% SNmono ha "remark" ma non "osservazione" in italiano.
\newtheorem{ossinner}[teorema]{Osservazione}
\newenvironment{oss}[1][]{%
	\begin{ossinner}[#1]\renewcommand{\qedsymbol}{\ensuremath{\square}}%
		\pushQED{\qed}\ignorespaces
	}{%
		\popQED\end{ossinner}%
}
\newtheorem{osservazione}[teorema]{Osservazione}
\newtheorem*{nota}{Nota}

% ── Esempio ──────────────────────────────────────────────────────────────
% SNmono ha "example" ma non "esempio" in italiano.
\newtheorem{esempioinner}{Esempio}[chapter]
\newenvironment{esempio}[1][]{%
	\begin{esempioinner}[#1]\renewcommand{\qedsymbol}{\ensuremath{\blacksquare}}%
		\pushQED{\qed}\ignorespaces
	}{%
		\popQED\end{esempioinner}%
}

% ── Esercizio ────────────────────────────────────────────────────────────
% SNmono ha "exercise" ma non "esercizio" in italiano.
\newtheorem{esercizioinner}{Esercizio}[chapter]
\newenvironment{esercizio}[1][]{%
	\begin{esercizioinner}[#1]\renewcommand{\qedsymbol}{\ensuremath{\blacksquare}}%
		\pushQED{\qed}\ignorespaces
	}{%
		\popQED\end{esercizioinner}%
}

% ── Svolgimento ──────────────────────────────────────────────────────────
\newenvironment{svolgimento}{%
	\par\smallskip\noindent\textit{Svolgimento.}\space\ignorespaces
}{\par}

% ── Sotto-titolo e risultato ─────────────────────────────────────────────
\newcommand{\esottotit}[1]{\par\smallskip\noindent\textbf{#1}}
\newcommand{\risultato}{\par\smallskip\noindent\textit{Risultato.}\space}

% ── Sommario (alias di abstract) ─────────────────────────────────────────
\newenvironment{sommario}
{\par\addvspace{6pt}\noindent\textbf{Abstract}\\[4pt]\ignorespaces}
{\par\addvspace{6pt}}

% ── Proof in italiano ────────────────────────────────────────────────────
\addto\captionsitalian{\renewcommand{\proofname}{Dimostrazione}}